\chapter{Flu Vaccine FAQs}

\section*{Definition and Overview}
The influenza (flu) vaccine is a safe and effective way to protect against seasonal influenza, a serious viral illness that hospitalises and kills thousands of Australians each year. The vaccine is updated every year to match the circulating strains, which is why it needs to be given annually. It is recommended for everyone aged 6 months and over, and it is free for groups at higher risk, including older adults, children under 5, pregnant women, Aboriginal and Torres Strait Islander people, and people with chronic conditions. This chapter answers the most common questions about the flu vaccine, including how it works, who should have it, side effects, and how it differs from COVID-19 vaccines.

\section*{Epidemiology}
Influenza causes substantial illness every year in Australia. In a typical season, it infects hundreds of thousands of people, causes tens of thousands of GP visits, and is associated with thousands of hospitalisations and hundreds of deaths, mostly among older adults and people with chronic disease. Vaccination coverage is around 50--60\% of the population, with higher rates in the groups for whom it is free. When coverage is high and the vaccine is a good match, the burden of disease falls markedly. In years with a good match, the vaccine reduces the risk of influenza illness by around 40--60\% and reduces hospitalisation and death even more.

\section*{Aetiology and Pathophysiology}
Understanding how the vaccine works answers many common questions.
\begin{itemize}
    \item \textbf{The virus:} influenza A and B viruses cause seasonal flu, and they change (mutate) constantly, which is why the vaccine is updated yearly
    \item \textbf{How the vaccine works:} it trains the immune system to recognise the strains expected that season, producing antibodies that fight infection
    \item \textbf{Inactivated vaccine:} the standard flu vaccine contains killed or weakened virus components and cannot cause influenza
    \item \textbf{Annual revaccination:} immunity fades over months, and the circulating strains change, so a yearly dose is needed
    \item \textbf{Protection timing:} it takes about 2 weeks for full protection to develop
    \item \textbf{What it protects against:} influenza, not colds or COVID-19; it is not 100\% effective but reduces severe illness, hospitalisation, and death
    \item \textbf{Herd immunity:} high vaccination rates protect vulnerable people who cannot be vaccinated
    \item \textbf{Safety:} the vaccine has an excellent safety record, with serious side effects being extremely rare
\end{itemize}

\section*{Clinical Features}
The features of vaccination are the practical details people ask about.
\begin{itemize}
    \item A single injection, usually in the upper arm, given from around April each year before the winter season
    \item Protection begins about 2 weeks after vaccination and lasts for the season
    \item Common side effects: a sore arm, mild headache, low-grade fever, and tiredness for a day or two
    \item These side effects reflect the immune response, not influenza
    \item The vaccine cannot give you the flu, because it contains no live virus
    \item Children may have a nasal spray version in some programs
    \item It can be given at the same time as a COVID-19 vaccine
    \item It is safe in pregnancy, and protects both the mother and the baby after birth
\end{itemize}

\section*{Differential Diagnosis}
Deciding about vaccination involves weighing common concerns.
\begin{itemize}
    \item The belief that the vaccine causes the flu, which it does not
    \item Concern about egg allergy: most people with egg allergy, including severe allergy, can now have the flu vaccine safely
    \item Confusion between flu, colds, and COVID-19, which are different illnesses
    \item The question of whether vaccination is needed every year, which it is
    \item Concern about side effects versus the risks of influenza, which are far greater
    \item Guillain--Barré syndrome, a very rare association that must be weighed against the much larger risk of flu complications
\end{itemize}

\section*{When to Seek Medical Care}
Most people do not need medical care for vaccination side effects. Seek advice if symptoms are severe, or if a fever and cough develop within days of vaccination and might actually be influenza.

\textbf{Call 000 (or local emergency services) for:}
\begin{itemize}
    \item Signs of a severe allergic reaction within minutes of vaccination: difficulty breathing, swelling of the face or throat, or collapse
    \item Severe symptoms, weakness, or difficulty moving after vaccination
    \item Any life-threatening emergency
\end{itemize}

\section*{Diagnosis and Investigations}
\begin{itemize}
    \item \textbf{Eligibility:} the vaccine is recommended for everyone from 6 months of age and is free for at-risk groups
    \item \textbf{Timing:} vaccination is best given from April onwards, before the peak of the season
    \item \textbf{Contraindications:} previous severe allergic reaction to a flu vaccine is the main reason to avoid it; mild illness is not a reason to delay
    \item \textbf{Co-administration:} the flu vaccine can be given with COVID-19 and other vaccines
    \item \textbf{Record keeping:} vaccination is recorded on the Australian Immunisation Register
\end{itemize}

\section*{Management}
\begin{itemize}
    \item \textbf{Where to get it:} GPs, pharmacies, and community immunisation clinics offer the vaccine from around April
    \item \textbf{Who gets it free:} people aged 65 and over, children aged 6 months to under 5, pregnant women, Aboriginal and Torres Strait Islander people, people with medical conditions such as heart, lung, kidney, and liver disease and diabetes, and other eligible groups
    \item \textbf{Children:} children aged 6 months to 5 years are at higher risk of severe flu and are eligible for a free vaccine
    \item \textbf{Pregnancy:} vaccination is recommended at any stage of pregnancy and protects the baby for the first months of life
    \item \textbf{Older adults:} a higher-dose or adjuvanted vaccine is used for people over 65 for stronger protection
    \item \textbf{Workplace programs:} many workplaces offer free or subsidised flu vaccination
    \item \textbf{After vaccination:} report serious side effects, and remember that other measures, including hand hygiene and staying home when unwell, still matter
\end{itemize}

\section*{Complications}
\begin{itemize}
    \item Influenza itself: pneumonia, bronchitis, worsening of chronic disease, hospitalisation, and death in severe cases
    \item Vaccine side effects, which are minor and short-lived in most people
    \item Very rare serious reactions, including severe allergy and Guillain--Barré syndrome
    \item Reduced protection in seasons when the vaccine is a poor match to circulating strains
    \item Missed protection when people do not get vaccinated
\end{itemize}

\section*{Prognosis and Outcomes}
The flu vaccine is the single best protection against seasonal influenza. In people who are vaccinated, the risk of influenza illness falls by roughly half in well-matched seasons, and the risks of hospitalisation, intensive care, and death are reduced substantially. Even when the match is imperfect, the vaccine reduces the severity of illness and protects the most vulnerable. Yearly vaccination is safe, quick, and widely available, and its benefits vastly outweigh its risks for nearly everyone.

\section*{Prevention and Self-Care}
\begin{itemize}
    \item Book the flu vaccine from April each year, before the season peaks
    \item Check if you are in a free-eligibility group
    \item Get vaccinated if pregnant, and protect children from 6 months of age
    \item Combine the flu vaccine with good hand hygiene and staying home when unwell
    \item Keep the vaccination record up to date
    \item Ask your pharmacist or GP about the vaccine and any concerns
    \item Do not let the common side effects of a sore arm or mild fever deter you; they are brief and minor
\end{itemize}

\section*{Sources and Further Reading}
The following reputable sources provide further information for healthcare students and patients seeking reliable, current advice about the flu vaccine.
\begin{itemize}
    \item Australian Government Department of Health and Aged Care. \textit{Influenza vaccination.} https://www.health.gov.au/
    \item Healthdirect Australia. \textit{Flu vaccine.} https://www.healthdirect.gov.au/flu-vaccine
    \item National Centre for Immunisation Research and Surveillance. \textit{Flu vaccine FAQs.} https://www.ncirs.org.au/
    \item The Royal Australian College of General Practitioners. \textit{Flu vaccination.} https://www.racgp.org.au/
    \item NHS. \textit{Flu vaccine.} https://www.nhs.uk/conditions/vaccinations/flu-influenza-vaccine/
    \item World Health Organization. \textit{Influenza vaccination.} https://www.who.int/
\end{itemize}
